% GENERATED FILE. Edit canonical lessons, then run npm run generate:books.
% canonical-source-hash: fnv1a64:507da1017fd24703
% canonical-lessons: PT-C309-nascimento, PT-C309-parto, PT-C309-gravida, PT-C309-gravidez, PT-C309-recem-nascido

\chapter{Nasceu o bebé}
\label{ch:pt-nasceu-o-bebe}
\hlchaptermodality{309}

\begin{chapteropening}
\textbf{By the end of this chapter:} \emph{I can talk about a birth: pregnancy, the delivery and the newborn baby.}

Complete the last lesson of chapter 309: I can talk about a birth: pregnancy, the delivery and the newborn baby.
\end{chapteropening}

\section[o nascimento]{o nascimento — the birth}

\label{lesson:PT-C309-nascimento}



\begin{quote}
Before the new one: say the Portuguese for to faint, then the Portuguese for to twist.
\end{quote}

\subsection*{You'll want to know: o nascimento}
\textbf{o nascimento} — \textquotedblleft{}the birth\textquotedblright{}.

\begin{quote}
O nascimento da minha filha mudou a minha vida.
\end{quote}

\begin{quote}
\emph{My daughter’s birth changed my life.}
\end{quote}

\emph{Data de nascimento} — date of birth.

\begin{cousinweb}
\emph{nascimento} is built on \emph{nascer} \textquotedblleft{}to be born\textquotedblright{}, which goes back to Latin \emph{nāscī} \textquotedblleft{}to be born\textquotedblright{} (Lewis \& Short), with the ending \emph{-mento}.
\end{cousinweb}

\begin{grammarlens}[title={What you've built}]
One more word for telling, explaining and describing at B1.
\end{grammarlens}

\subsection*{Your turn}
\begin{itemize}
\raggedright
  \item \emph{Say it:} \emph{o nascimento}
  \item \emph{Say it:} \emph{o nascimento}, once more
  \item \emph{Say it:} a sentence of your own with \emph{o nascimento}
  \item \emph{Recall:} say the Portuguese for to faint, then the Portuguese for to twist, then say \emph{o nascimento} again
\end{itemize}

\begin{tcolorbox}[breakable,colback=teal!4,colframe=teal!35!black,title={Before you move on}]
What does o nascimento mean? (\textquotedblleft{}The birth\textquotedblright{}.) Say it once more.
\end{tcolorbox}
\section[o parto]{o parto — the delivery, childbirth}

\label{lesson:PT-C309-parto}



\begin{quote}
Before the new one: say the Portuguese for to twist, then the Portuguese for the birth.
\end{quote}

\subsection*{You'll want to know: o parto}
\textbf{o parto} — \textquotedblleft{}the delivery, childbirth\textquotedblright{}.

\begin{quote}
O parto correu bem.
\end{quote}

\begin{quote}
\emph{The delivery went well.}
\end{quote}

\emph{Estar em trabalho de parto} — to be in labour.

\begin{cousinweb}
\emph{parto} is inherited from Latin \emph{partus} \textquotedblleft{}a birth\textquotedblright{} (EtymDB), from \emph{parere} \textquotedblleft{}to give birth\textquotedblright{}.
\end{cousinweb}

\begin{grammarlens}[title={What you've built}]
One more word for telling, explaining and describing at B1.
\end{grammarlens}

\subsection*{Your turn}
\begin{itemize}
\raggedright
  \item \emph{Say it:} \emph{o parto}
  \item \emph{Say it:} \emph{o parto}, once more
  \item \emph{Say it:} a sentence of your own with \emph{o parto}
  \item \emph{Recall:} say the Portuguese for to twist, then the Portuguese for the birth, then say \emph{o parto} again
\end{itemize}

\begin{tcolorbox}[breakable,colback=teal!4,colframe=teal!35!black,title={Before you move on}]
What does o parto mean? (\textquotedblleft{}The delivery, childbirth\textquotedblright{}.) Say it once more.
\end{tcolorbox}
\section[grávida]{grávida — pregnant}

\label{lesson:PT-C309-gravida}



\begin{quote}
Before the new one: say the Portuguese for the birth, then the Portuguese for the delivery.
\end{quote}

\subsection*{You'll want to know: grávida}
\textbf{grávida} — \textquotedblleft{}pregnant\textquotedblright{}.

\begin{quote}
A minha irmã está grávida de seis meses.
\end{quote}

\begin{quote}
\emph{My sister is six months pregnant.}
\end{quote}

\emph{Estar grávida de} + the months.

\begin{cousinweb}
\emph{grávida} was borrowed from Latin \emph{gravida} \textquotedblleft{}pregnant\textquotedblright{} (EtymDB), from \emph{gravis} \textquotedblleft{}heavy\textquotedblright{}.
\end{cousinweb}

\begin{grammarlens}[title={What you've built}]
One more word for telling, explaining and describing at B1.
\end{grammarlens}

\subsection*{Your turn}
\begin{itemize}
\raggedright
  \item \emph{Say it:} \emph{grávida}
  \item \emph{Say it:} \emph{grávida}, once more
  \item \emph{Say it:} a sentence of your own with \emph{grávida}
  \item \emph{Recall:} say the Portuguese for the birth, then the Portuguese for the delivery, then say \emph{grávida} again
\end{itemize}

\begin{tcolorbox}[breakable,colback=teal!4,colframe=teal!35!black,title={Before you move on}]
What does grávida mean? (\textquotedblleft{}Pregnant\textquotedblright{}.) Say it once more.
\end{tcolorbox}
\section[a gravidez]{a gravidez — pregnancy}

\label{lesson:PT-C309-gravidez}



\begin{quote}
Before the new one: say the Portuguese for the delivery, then the Portuguese for pregnant.
\end{quote}

\subsection*{You'll want to know: a gravidez}
\textbf{a gravidez} — \textquotedblleft{}pregnancy\textquotedblright{}.

\begin{quote}
Durante a gravidez, ela continuou a trabalhar.
\end{quote}

\begin{quote}
\emph{During her pregnancy she went on working.}
\end{quote}

Built on \emph{grávida}.

\begin{cousinweb}
\emph{gravidez} is built on \emph{grávida} \textquotedblleft{}pregnant\textquotedblright{}, from Latin \emph{gravidus} \textquotedblleft{}laden, heavy\textquotedblright{}, from \emph{gravis} \textquotedblleft{}heavy\textquotedblright{} (Lewis \& Short).
\end{cousinweb}

\begin{grammarlens}[title={What you've built}]
One more word for telling, explaining and describing at B1.
\end{grammarlens}

\subsection*{Your turn}
\begin{itemize}
\raggedright
  \item \emph{Say it:} \emph{a gravidez}
  \item \emph{Say it:} \emph{a gravidez}, once more
  \item \emph{Say it:} a sentence of your own with \emph{a gravidez}
  \item \emph{Recall:} say the Portuguese for the delivery, then the Portuguese for pregnant, then say \emph{a gravidez} again
\end{itemize}

\begin{tcolorbox}[breakable,colback=teal!4,colframe=teal!35!black,title={Before you move on}]
What does a gravidez mean? (\textquotedblleft{}Pregnancy\textquotedblright{}.) Say it once more.
\end{tcolorbox}
\section[o recém-nascido]{o recém-nascido — the newborn baby}

\label{lesson:PT-C309-recem-nascido}



\begin{quote}
Before the new one: say the Portuguese for pregnant, then the Portuguese for pregnancy.
\end{quote}

\subsection*{You'll want to know: o recém-nascido}
\textbf{o recém-nascido} — \textquotedblleft{}the newborn baby\textquotedblright{}.

\begin{quote}
O recém-nascido dorme quase o dia todo.
\end{quote}

\begin{quote}
\emph{The newborn sleeps almost all day.}
\end{quote}

\emph{Recém-} means ‘newly’: \emph{recém-casados}.

\begin{cousinweb}
\emph{recém-nascido} joins \emph{recém}, a short form of \emph{recente}, from Latin \emph{recēns} \textquotedblleft{}fresh, new\textquotedblright{}, and \emph{nascido} \textquotedblleft{}born\textquotedblright{}, from Latin \emph{nāscī} (Lewis \& Short).
\end{cousinweb}

\begin{grammarlens}[title={What you've built}]
One more word for telling, explaining and describing at B1.
\end{grammarlens}

\subsection*{Your turn}
\begin{itemize}
\raggedright
  \item \emph{Say it:} \emph{o recém-nascido}
  \item \emph{Say it:} \emph{o recém-nascido}, once more
  \item \emph{Say it:} a sentence of your own with \emph{o recém-nascido}
  \item \emph{Recall:} say the Portuguese for pregnant, then the Portuguese for pregnancy, then say \emph{o recém-nascido} again
\end{itemize}

\begin{tcolorbox}[breakable,colback=teal!4,colframe=teal!35!black,title={Before you move on}]
What does o recém-nascido mean? (\textquotedblleft{}The newborn baby\textquotedblright{}.) Say it once more.
\end{tcolorbox}
